\chapter{Conclusion}
\label{chap:conclusion}

This work developed a hot-started nullspace active-set method for the quadratic master
problems of regularized column generation, for multi-block linear programs whose linking
constraints may be equalities or inequalities of either direction.

\section{Summary}

The starting point was a \textbf{unified formulation} (\cref{chap:formulation}): surplus
variables bring every linking row into equality form, and the sign conditions of the
inequality rows become ordinary dual constraints of the surplus columns. Proximal
regularization of the dual then leads to a convex quadratic master problem
(\cref{chap:regularization}). Because all linking rows are equalities, its auxiliary variable
can be eliminated for all rows, and the result has only the convexity constraints and
nonnegativity bounds. Two identities underlie everything that follows: the duals of the
linking rows are an explicit affine function of the primal iterate, and the gradient of the
quadratic objective consists of the reduced costs of the columns and the duals of the
inequality rows.

\medskip

For the active-set method (\cref{chap:qpas}), the constraint matrix of a working set admits a
\textbf{reordering into five sets}: representatives, working-set columns, nullspace columns,
active surpluses and free surpluses. The surpluses violate the column generation property of
the convexity matrix, but this only adds a group with identity structure (active surpluses)
and a group of empty columns (free surpluses). The
reordering yields an explicit nullspace basis of vectors $e_j - e_{\rep(j)}$ and unit vectors,
which is never stored (\cref{chap:nullspace}). The reduced KKT system \textbf{decouples}: the
surplus part of the direction is explicit, and only the \textbf{effective reduced Hessian}
$\Heff = (\Aeff Z_x)^\top(\Aeff Z_x)$ has to be factorized. Here $\Aeff$ holds the equality rows
and the inequality rows with active surplus. The column part of the method is thus the method
for equality rows, applied to a linking matrix whose rows change with the surplus working set.
The Lagrange multipliers are the reduced costs and the dual signs, so the optimality test of
the master problem and the pricing test are the same test.

\medskip

Every change of the working set is a row addition, a row deletion or a rank-one modification
of an $LDL^\top$ factorization of $\Heff$ (\cref{chap:updates}). A new column or a new
stability center leaves the factorization unchanged, which gives a hot start of every master
problem. Singular reduced Hessians are resolved by an exact zero-curvature step, which removes
the dependency and keeps every working-set variable except the released one at its bound, for
released columns and released surpluses alike (\cref{chap:singular}).

\medskip

The factorization was realized in two ways. A sparse factorization of fixed dimension
$\nmax$, updated with CHOLMOD, accumulates structural zeros, which periodic refactorization in
increasing index order controls (\cref{chap:updates}). Since $\Heff$ is small and dense, it can
instead be factorized densely in the order in which its indices became active
(\cref{chap:dense}). A released column is then appended, a removed index is deleted by a
rank-one update of the positions behind it, and the surplus operations are rank-one updates
and downdates by the method of Gill, Golub, Murray and Saunders. This factor never contains a
stale entry and needs no periodic refactorization; its new pivot is exactly the Schur
complement of a released column, its downdates fail exactly when the downdated matrix is
singular, and removals never decrease a pivot.

\medskip

On generated crew scheduling instances (\cref{chap:results}), the method reproduces the optimum
of the full linear program and certifies it by the Lagrangian bound, with both factorizations.
With the dense factorization, its master problems cost between $0.003$ and $13$\,ms on
average, almost three orders of magnitude less than a general-purpose QP solver, and fast
enough to make regularized column generation much faster than classical column generation with
a warm-started LP master: 3.0\,s instead of 663\,s on the largest instance on which both were
run. The dense factorization is two to four times faster than the sparse one, with a speedup
that grows with the number of linking rows, needs an order of magnitude less memory, and is
just as accurate. The complete column generation run is faster than solving the full linear
program directly from the second instance on, by a factor of three to four on the larger
instances, and on the largest ones the direct solve did not finish in the time in which the
active-set method was done. Inequality rows need no new kind of linear algebra, keep the time
per master problem of the same order, and lead to shorter runs.


\section{Outlook}

The experiments point to four directions for further work.
\begin{itemize}
    \item \textbf{A cheaper multiplier test.} With the dense factorization, the linear algebra
          is no longer the dominant cost. The multiplier test, which computes the reduced cost of
          every working-set column at every subspace minimizer, now costs as much as the
          direction computation. Since a step changes $\pi$ only by a multiple of $\Abar p$, the
          reduced costs could be updated incrementally, or the test could be restricted to a
          candidate list as in partial pricing.
    \item \textbf{Adaptive stabilization in the proximal phase.} The long tail of center-only
          iterations is a property of the proximal point method with a fixed $\delta$
          (\cref{prop:center-only}). Since $\delta$ does not enter the factorization, it can be
          decreased in this phase at no cost to the linear algebra.
    \item \textbf{Range rows.} Linking rows with two finite bounds need surpluses with upper
          bounds, which fit into the same framework as an additional working-set state.
    \item \textbf{Inexact master solves.} Since the multipliers are the reduced costs, pricing
          could be called at any subspace minimizer instead of only at the optimum of the
          master problem. This would turn column generation into a single active-set method
          over all columns.
\end{itemize}
